\documentclass[11pt]{amsart}
\usepackage{geometry}                % See geometry.pdf to learn the layout options. There are lots.
\geometry{letterpaper}                   % ... or a4paper or a5paper or ... 
%\geometry{landscape}                % Activate for for rotated page geometry
%\usepackage[parfill]{parskip}    % Activate to begin paragraphs with an empty line rather than an indent
\usepackage{graphicx}
\usepackage{amssymb}
\usepackage{epstopdf}
\DeclareGraphicsRule{.tif}{png}{.png}{`convert #1 `dirname #1`/`basename #1 .tif`.png}

\title{PS 4.20}
%\date{}                                           % Activate to display a given date or no date

\begin{document}
\maketitle
%\section{}
%\subsection{}

\noindent Sean A(M x N) y B(N) arreglos de dos y una dimensión, respectivamente. Cons­ truya un diagrama de flujo que asigne valores a A, a partir de B teniendo en cuen­ ta los siguientes criterios:

\[
\begin{aligned}
a) &\quad A[i,j] = B[i]  \quad Si \quad i \leq j \\
b) &\quad A[i,j] = 0 \quad Si \quad i > j 
\end{aligned}
\]
\newline


\noindent Datos: A[1..M, 1..N], B[1..N]\quad  (1 $\leq$ M $\leq$ 30, 1 $\leq$ N $\leq$ 20)
\newline

\noindent Donde:
\newline

\noindent A es un arreglo bidimensional de tipo real de M renglones y N co­lumnas.

\noindent B es un arreglo unidimensional de N elementos de tipo real.


\end{document}  